% Source: https://artofproblemsolving.com/community/c3h1402962_lines_and_more_lines
% Stable identity: aops-topic-1402962-q01
\textbf{Upper bound.} Two distinct lines meet in at most one point. Every intersection point lies on at least two of the lines, so it is the meeting point of at least one pair of lines. Assigning to each intersection point one such pair gives different pairs for different points (a pair of lines meets in only one point). Hence the number of intersection points is at most the number of pairs of lines,
\[
\dbinom{n}{2}=\dfrac{n(n-1)}{2}.
\]
\textbf{The bound is attained} when no two lines are parallel and no three pass through a common point, because then each pair gives its own point. For an explicit example take the lines
\[
\ell_k:\ y=kx+k^2\qquad(k=1,2,\dots,n).
\]
Lines $\ell_k$ and $\ell_m$ ($k\ne m$) have different slopes, and solving $kx+k^2=mx+m^2$ gives $x=-(k+m)$, $y=-km$. The point $\bigl(-(k+m),-km\bigr)$ determines the sum and the product of $k$ and $m$, hence the pair $\{k,m\}$ itself (they are the roots of $t^2-(k+m)t+km$). So different pairs give different points, and there are exactly $\dfrac{n(n-1)}{2}$ intersection points.
